% Source: https://artofproblemsolving.com/community/c3h1247885_geometry
% Stable identity: aops-topic-1247885-q01
The volume of a tank of radius $r$ is $V(r)=10\pi r^2+\dfrac{4}{3}\pi r^3$ (the two hemispheres form a sphere).

\begin{asy}
usepackage("amsmath");
size(5cm);
defaultpen(fontsize(10pt));
real r=3,a=10;
draw((0,r)--(a,r)); draw((0,-r)--(a,-r));
draw(arc((0,0),r,90,270)); draw(arc((a,0),r,-90,90));
draw((0,0)--(0,r),dashed); label("$r$",(0,r/2),E);
draw((0,-r-1)--(a,-r-1),Arrows); label("$10$",(a/2,-r-1),S);
label("\mbox{}",(min(currentpicture,user=true)+max(currentpicture,user=true))/2,black); // keeps the text after the figure black when LaTeX typesets the labels inline
\end{asy}

The standard tank has $r=3$: $V=90\pi+36\pi=126\pi$. We need $10\pi r^2+\dfrac{4}{3}\pi r^3=252\pi$, i.e.
\[2r^3+15r^2-378=0.\]
The left side increases for $r>0$; it is $-10$ at $r=4$ and about $+0.9$ at $r=4.05$, and Newton's method or bisection gives $r\approx\boxed{4.05}$ feet (more precisely $4.046$).
